\documentclass[10pt,a4paper]{article}
\usepackage{resume-style}
\usepackage{xeCJK}
\setmainfont{Avenir Next}
\setCJKmainfont[
  Path=/System/Library/Fonts/,
  UprightFont={STHeiti Light.ttc},
  BoldFont={STHeiti Medium.ttc}
]{STHeiti Light.ttc}
\urlstyle{same}
\hypersetup{
  pdftitle={王泽森 - 中文简历},
  pdfauthor={王泽森},
  pdfsubject={机器人系统、VLA 研究与 VLA 后训练}
}

\begin{document}

\cvname{王泽森 \fontsize{13}{15}\selectfont Vincent Wang}
\cvtagline{机器人系统 \textbar{} 视觉-语言-动作模型 \textbar{} VLA 后训练}
\cvcontact{美国宾夕法尼亚州费城 \enspace\textbar\enspace +1 267-257-9120 \enspace\textbar\enspace
  \href{mailto:wang2003@engineering.upenn.edu}{wang2003@engineering.upenn.edu} \enspace\textbar\enspace
  \href{https://www.linkedin.com/in/sansenpai/}{LinkedIn} \enspace\textbar\enspace
  \href{https://github.com/oldkingzz}{GitHub} \enspace\textbar\enspace
  \href{https://oldkingzz.github.io/profile/}{个人网站}}

\cvsection{教育经历}
\cventry{宾夕法尼亚大学 University of Pennsylvania}{预计 2027 年 5 月毕业}{机械工程与应用力学硕士 - 机电与机器人方向}
\begin{cvitems}
  \item 相关课程：高等动力学、机电系统、工程数学、操作系统。
\end{cvitems}

\cventry{华东理工大学 East China University of Science and Technology}{2025 年 6 月}{智能制造工程学士 \textbar{} GPA：3.6/4.0 \textbar{} 专业前 15\%}
\vspace{1.4pt}

\cvsection{实习经历}
\cventry{Synthoid.ai \textbar{} VLA 实习研究员}{2026.05.18 - 2026.08.21}{真机基础设施、VLA 训练与独立策略蒸馏研究}
\begin{cvitems}
  \item 独立完成 7 自由度工业机械臂的机器人端 SDK 控制与 VLA adapter，实现笛卡尔末端位姿控制、分层控制频率、轨迹插值和自动恢复；交付稳定的人类遥操作链路。
  \item 主导灵巧手、头部/腕部相机与 LeRobot 数据格式下的遥操作和真机数据采集；完成多模态时间戳同步，并基于动作一至三阶导数筛查异常轨迹。
  \item 深度参与公司自研 VLA 模型的前期训练，负责训练实现与配置、分布式运行、checkpoint 管理、初始消融和故障诊断；定位 gripper 维度输出契约不一致问题。
  \item 独立定义并实现面向 flow-matching VLA 蒸馏的 Path-OPD。在 outer on-policy occupancy、Teacher 查询、训练与评测预算匹配的条件下，closed-loop success 明显优于 endpoint DAgger；正在准备 ICLR 2027 投稿。
\end{cvitems}

\cventry{上海无极科技 \textbar{} 机器人控制实习生}{2024 年 7 月 - 9 月}{灵巧手遥操作与底层电机控制}
\begin{cvitems}
  \item 在 MATLAB/Simulink 中将不稳定的阻抗控制环重构为带力矩反馈的导纳控制器，用于灵巧手遥操作。
  \item 开发 Linux C++ 上位机，实现约 30 个电机的可靠 CAN 通信与底层控制。
\end{cvitems}

\cvsection{代表性研究与领导经历}
\cventry{华东理工大学 Robocon 战队 \textbar{} 联合创始人、控制负责人}{2022 年 9 月 - 2024 年 7 月}{自主机器人系统与跨学科技术领导}
\begin{cvitems}
  \item 从零搭建战队首套 STM32 运动控制系统，并将 Point-LIO、YOLO（嵌入式约 20 FPS）、MoveIt 与底盘控制集成为自主机器人系统。
  \item 带领 20 人跨学科团队，负责系统架构与技术选型，并组织 60 名新成员的技术培训。
\end{cvitems}

\cvcompactentry{MDATS - 面向机械故障检测的多域自适应分类方法}{IEEE MEAE 2024}
\begin{cvitems}
  \item 第一作者；搭建 PyTorch 研究流程，目标域准确率达到 88.27\%，较 baseline 提升 40.6 个百分点。
\end{cvitems}

\cvsection{技术能力}
\cvskill{编程与系统}{Python、C/C++、Rust、Linux、Git、CMake、ROS 2、CAN}
\cvskill{机器人与控制}{笛卡尔控制、遥操作、阻抗/导纳控制、STM32/ESP32、FreeRTOS、MoveIt}
\cvskill{学习与研究}{PyTorch、分布式训练、VLA 后训练、flow matching、policy distillation、LeRobot、对照实验设计}

\end{document}
